\documentclass[11pt]{article}
\usepackage[margin=1in]{geometry}
\usepackage[utf8]{inputenc}
\usepackage[T1]{fontenc}
\usepackage{lmodern}
\usepackage{amsmath,amssymb}
\usepackage{graphicx}
\usepackage{booktabs}
\usepackage{xcolor}
\usepackage[hidelinks]{hyperref}
\usepackage{url}
\usepackage{caption}
\captionsetup{font=small,labelfont=bf}

\newcommand{\memprint}{MemPrint}
\newcommand{\secref}[1]{Section~\ref{#1}}
\newcommand{\figref}[1]{Figure~\ref{#1}}
\newcommand{\tabref}[1]{Table~\ref{#1}}
\newcommand{\term}[1]{\emph{#1}}
\graphicspath{{figures_motivation/}}

\title{What Does a Scheduler Lose Without an Exact Memory Footprint?\\
Existing Estimators and Provisioning Policies on Disaggregated Memory}
\author{Nafis Mustakin\\University of California, Riverside}
\date{October 2026}

\begin{document}
\maketitle

\begin{abstract}
A scheduler for a cluster with disaggregated memory can pack jobs only as tightly as its knowledge of each job's
memory footprint allows. We measure how much an exact footprint is worth to such a scheduler, and how much of that
worth the existing sources of footprint knowledge recover. We trace 25 programs at three inputs each, HPC kernels,
interpreters, a compiler and two LLM inference engines among them, with a full instrumentation trace as the truth and
with Massif, DHAT, heaptrack, the gperftools heap profiler and sampled RSS, and we replay the measured footprints
through DMSched-Sim, a simulator of a cluster with local DRAM and a shared memory pool in which only the footprint
oracle changes. We also evaluate three published provisioning policies, Pond, the dynamic provisioning of Zacarias et
al.\ and KS+, and replay jobs of the SURF Lisa trace with their measured memory. At 80\% offered load the exact
footprint curve gives a mean response time 26.4$\times$ shorter than the user's request and 5.3$\times$ shorter than
the exact peak, and it needs 5.0\,TB of memory for a service level that the exact peak meets with 6.0\,TB and the
request with 11.0\,TB. No existing source recovers this gain. The heap profilers report the right peak for most
programs but read a compiler and an LLM engine at 0.14 to 0.17 of their peaks, because they see only the heap, so no
margin makes them safe without over-reserving every other job. Pond holds the whole request and is slower than the
request itself. Dynamic provisioning comes close to the exact peak only by reading the coming minutes of each job's
use from its trace; a version that looks back, as a running system must, falls behind the exact peak. KS+'s
regressions on input size fail when the job's input is far from the ones it has seen. On the Lisa trace the exact curve
completes 18 to 31\% more jobs than the exact peak once memory is scarce. The exact curve itself costs 227 times the
program's run to obtain. We derive from these results the requirements a footprint estimator must meet to be useful
to a scheduler.
\end{abstract}

\section{Introduction}
\label{sec:intro}

A batch scheduler places a job only after it decides how much memory to reserve for it, and on a cluster with
disaggregated memory that decision also sets how much of the shared pool the job takes from other jobs. Disaggregated
memory attaches a pool of memory, over CXL or RDMA, that every node can use in addition to its local
DRAM~\cite{Li2023Pond,Zacarias2021Throughput}. The pool lets a cluster provision less local memory per node, but only if
the scheduler can pack jobs tightly, and it can pack them only as tightly as its knowledge of each job's memory allows.
That knowledge is weak today. Peng et al.\ found that over 90\% of LLNL jobs use under 15\% of a node's
memory~\cite{Peng2020Underutilization}, and Zacarias et al.\ showed that a uniform 60\% overestimate of memory demand cuts
the throughput of a disaggregated system by about 25\% and raises the mean response time about five
times~\cite{Zacarias2021Accurate}.

Two kinds of knowledge are possible. A \term{peak} is one number per job, held for the whole run. A \term{footprint
curve} gives the memory the job needs at every point of its run, so a scheduler can let two jobs whose peaks do not
overlap share the same memory. The usual sources of either are the user's request, a heap profiler such as
Massif~\cite{ValgrindMassif}, DHAT~\cite{ValgrindDHAT}, heaptrack~\cite{heaptrack} or the gperftools heap
profiler~\cite{gperftools}, the resident set size (RSS) of an earlier run, and three kinds of published policy: a
learned predictor of a job's untouched memory (Pond~\cite{Li2023Pond}), a monitor that resizes a running job's allocation
(dynamic provisioning~\cite{Zacarias2023Dynamic}), and a model of a task's memory over time learned from its earlier
executions (KS+~\cite{KSplus2024}).

However, it is unclear how much a scheduler loses with these sources against an exact footprint, and why. The
simulators that study scheduling on disaggregated memory either model one run in cycle detail with no
scheduler~\cite{Puri2024DRackSim,CXLClusterSim,CXLDMSim,Yang2023CXLMemSim}, or treat a job's memory as one declared peak
with a synthetic error~\cite{Zacarias2021Accurate,Zacarias2021Throughput}. None feeds the scheduler the footprints that
real profilers report, and none compares them with an exact footprint curve. Therefore, we measure every source of
footprint knowledge listed above on the same programs and replay the results through one scheduler, with everything but
the footprint oracle held fixed.

In this work, we make the following contributions:
\begin{itemize}
  \item We build a corpus of footprint curves for 25 programs at three inputs each, with a full instrumentation trace
    as the truth, the curves and peaks of four heap profilers and sampled RSS, and what each method costs
    (\secref{sec:corpus}).
  \item We build DMSched-Sim, a cluster simulator with local memory and a shared pool, an EASY-backfilling scheduler
    that admits a job only if its whole reservation curve fits, a remote-memory slowdown model, OOM kills, and
    implementations of Pond, dynamic provisioning and KS+ checked against their authors' code where it is published
    (\secref{sec:sim}).
  \item We show how much of an exact footprint's gain comes from knowing the curve rather than the peak, how much each
    existing source recovers, why each falls short, and how much memory each needs for the same service
    (\secref{sec:results}), on a synthetic stream, on a stream of LLM jobs and on the SURF Lisa trace.
  \item We derive the requirements a footprint estimator must meet to be useful to a scheduler (\secref{sec:req}).
\end{itemize}

\section{Background and related work}
\label{sec:related}

\paragraph{Disaggregated memory simulators.} DRackSim~\cite{Puri2024DRackSim}, CXL-ClusterSim~\cite{CXLClusterSim},
CXL-DMSim~\cite{CXLDMSim} and CXLMemSim~\cite{Yang2023CXLMemSim} model the timing of accesses to pooled memory for one
run, in cycle detail or from a trace, and none of them has a job scheduler. The BSC Slurm simulator extended by Zacarias
et al.\ schedules real job traces on a disaggregated system~\cite{Zacarias2021Throughput}, but it represents a job's
memory as one peak, and its study of accuracy multiplies that peak by a synthetic overestimate~\cite{Zacarias2021Accurate}.
Boezennec et al.\ study the allocation of disaggregated memory to jobs with piecewise-constant memory
phases~\cite{Boezennec2024}, and CFM schedules jobs against per-job slowdown curves measured offline~\cite{Amaro2020CFM}.

\paragraph{Provisioning policies.} Pond predicts the untouched memory of virtual machines with a gradient-boosted model
from their customers' history, places that memory in a CXL pool, and moves a virtual machine to local memory when it
touches more than predicted~\cite{Li2023Pond}. The dynamic provisioning of Zacarias et al.\ resizes each job's allocation
every five minutes to its use, read from usage traces, and fails and restarts a job whose growth does not
fit~\cite{Zacarias2023Dynamic}. KS+ learns a workflow task's memory over time as a step function whose segments it
regresses on the input size of earlier executions, and reduces wasted GB-seconds by 51\% against peak-only
prediction~\cite{KSplus2024}. Pond and dynamic provisioning act on a job's allocation while it runs or from its owner's
history, and KS+ needs earlier executions of the same task.

\paragraph{Footprint estimators.} Massif samples the heap at up to 100 snapshots by default, on a time axis of
instructions or of milliseconds under Valgrind~\cite{ValgrindMassif}, and with \texttt{--pages-as-heap} it counts mapped
pages instead of heap blocks. DHAT reports the heap at its global peak (t-gmax) and at exit, and the time of the peak,
but no time series~\cite{ValgrindDHAT}. heaptrack records every allocation and can export a heap curve~\cite{heaptrack},
and the gperftools heap profiler writes a dump each time the heap grows by a set amount~\cite{gperftools}. All four count
requested heap bytes, so they miss memory a program maps itself and count blocks that are allocated but never touched.

\section{Simulator}
\label{sec:sim}

DMSched-Sim replays one stream of jobs once per footprint oracle, so that any difference between two runs comes from
what the scheduler knows about memory. It is a time-slotted simulator (one slot is one minute) written in Python.

\paragraph{System model.} The cluster has 32 nodes, each with 32 cores and 128\,GB of local DRAM, and one shared pool of
1\,TB (one sixth of all memory). A job's memory is placed in its node's free local memory first, and the rest goes to the
pool. A job on several nodes holds its footprint on each of them, its nodes share the pool, and it runs as long as its
slowest node.

\paragraph{Jobs.} A job has an arrival time, a core count, a base runtime (its runtime with all memory local), a true
footprint curve $a(p)$ over its progress $p \in [0, 1]$, and a reservation curve $r(p)$ set by the oracle. Curves are
defined over progress rather than time, so a job slowed by remote memory stretches its curve in wall-clock time.

\paragraph{Scheduler.} The scheduler is first come, first served with EASY backfilling~\cite{Lifka1995EASY} over a window
of 32 queued jobs. Each node keeps a timeline of committed local memory and cores, and the pool keeps one of committed
memory. A job starts on a node only if, at every slot from now to its projected end, the part of its reservation that
does not fit in the node's free local memory fits in the free pool. With a time-varying reservation, two jobs whose
peaks do not overlap can therefore share the same memory. When the head of the queue cannot start, the scheduler reserves
its earliest feasible start, and a later job may start early only if it fits beside that reservation. Among the nodes
that fit, the scheduler picks the one with the least pool use and then the least free local memory left over.

\paragraph{Remote-memory slowdown.} Following the linear models of CFM~\cite{Amaro2020CFM} and Boezennec et
al.~\cite{Boezennec2024}, a job's runtime is stretched by the share of its resident GB-minutes that lie above its local
reservation:
\begin{equation}
T = T_{\mathrm{base}}\,(1 + \beta f_{\mathrm{remote}}), \qquad
f_{\mathrm{remote}} = \frac{\sum_t \max\bigl(0,\, a(t) - \ell(t)\bigr)}{\sum_t a(t)},
\end{equation}
where $\ell(t)$ is the local part of the reservation and $\beta = 0.5$, so a job whose memory is entirely remote runs 1.5
times as long. The runtime and the placement are solved together by fixed-point iteration when the job is admitted.

\paragraph{OOM kills.} If the true footprint exceeds the reservation at any slot, the job is killed there, its work is
lost, and it is requeued at the head of the queue with its reservation multiplied by 1.5, as a batch system that raises a
job's memory limit after an out-of-memory failure would do. The policies that resize a running job handle growth as
\secref{sec:oracles} describes.

\paragraph{Metrics.} We report the mean response time (completion minus arrival), the mean bounded slowdown (response
over $\max(T_{\mathrm{base}}, 10\,\mathrm{min})$), the number of OOM kills, and the reserved and the resident memory as
shares of all memory, over the jobs that arrive between the 10th and the 90th percentile of arrival times.

\paragraph{Validation.} Unit tests check that two jobs with non-overlapping peaks share a node under their curves but run
one after the other under their peaks, that an under-reservation is killed where the footprint passes it, that
backfilling never delays the head of the queue, that a fully remote job runs $1 + \beta$ times as long, and the behaviour
of each published policy. The published policies are checked against their authors' code where it is published
(\secref{sec:oracles}), and Appendix~\ref{sec:sensitivity} reproduces the authors' own results.

\subsection{Job streams}
\label{sec:streams}

The main experiments use a controlled stream of 3{,}000 jobs in which memory is the contended resource. Each job has 1 to
16 cores, a lognormal base runtime with a median of 30 minutes clipped to 5 to 720 minutes, and a lognormal memory request
with a median of 100\,GB. What a job really uses is its request multiplied by a ratio $u \le 1$ taken from the LANL CM5
log of the Parallel Workloads Archive~\cite{PWA}, the archive log that records both requested and used memory (median
$u = 0.79$ over its 94{,}335 jobs with both fields). The shape of the job's footprint, and what every estimator reports
for it, come from one corpus program at its largest input (\secref{sec:corpus}), chosen at random, scaled so that the
truth's peak is the job's peak and every estimator keeps its own over- or under-estimate. Arrivals are spread over a span
set so that the \term{offered load}, the resident GB-minutes arriving per minute over the cluster's total memory, takes
the values 0.4 to 1.2. Each point is three seeds.

\paragraph{A real trace with measured memory.} The SURF Lisa trace~\cite{LisaTrace} joins 1.6 million Slurm jobs of a
343-node cluster with their nodes' Prometheus metrics, sampled every 30\,s from June to November 2022. A node's memory is a
job's only while the job has the node to itself, so we keep a job only if no other job shared any of its nodes at any
sample and at least 80\% of the samples its runtime implies are present: 88{,}767 jobs of at least five minutes, 1{,}429
of them on 2 to 22 nodes. A job's memory is its nodes' active memory above each node's idle level. We replay windows of
3{,}000 consecutive jobs in two ways. In the \term{measured} stream every job keeps its own curve, and only the oracles
that need nothing but the job's own memory apply. In the \term{hybrid} stream every job keeps the trace's arrival,
runtime, node count and measured peak, with a corpus program's footprint shape scaled to that peak, so that every
estimator applies. Lisa records no memory request, so a job requests Slurm's default on the simulated cluster, the node's
memory per core times the cores it uses, and at least its peak. Both streams run on a cluster of 32 nodes with 16 cores
each, as Lisa's, a pool of a quarter of the local memory, and local memory per node swept from Lisa's 103\,GB down to
4\,GB; the offered load is set against each cluster's own memory.

\subsection{Two truths}
\label{sec:truths}

A footprint can be counted as the bytes a program touches or as the pages a kernel holds for it, and a kernel charges and
OOM-kills a process for the latter. The two differ by page granularity, by memory an allocator keeps after a free, and by
the runtime's own pages. We run every experiment with two truths. With \term{bytes touched} as the truth, the scheduler
must cover the footprint a page-perfect system would need. With \term{RSS} as the truth, it must cover the resident set
sampled every 5\,ms.

\subsection{Oracles and margins}
\label{sec:oracles}

An oracle turns what is known about a job into its reservation curve. Curve oracles follow a measured curve, coarsened to
20 windows over progress, each holding its largest value and widened to its neighbours, because each method's clock
(instructions or milliseconds under Valgrind for Massif, seconds for the RSS) maps to wall-clock progress only
approximately. Peak oracles hold one value for the whole run. The \term{exact curve} and the \term{exact peak} are the
truth's own curve and peak, the bounds no estimator can pass. The estimators are Massif's curve (heap, or mapped pages
with \texttt{--pages-as-heap}) on its instruction clock and on its millisecond clock, heaptrack's curve, DHAT's peak and
two curves built from DHAT's three recorded points (\secref{sec:dhat}), gperftools' peak, the RSS peak of an earlier run
of the same input, the user's request, and an ML peak, the request times a 0.95-quantile gradient-boosted prediction of
$u$ from the request, core count, requested time and the user's and executable's history, trained on half of the CM5
users, as Pond's model predicts untouched memory~\cite{Li2023Pond}.

The policies that resize a running job are three. \term{Reactive resizing} holds the request for five minutes, then each
five minutes the past five minutes' largest use plus 10\% headroom, and grows the job from free memory when its use
exceeds its allocation. \term{Dynamic provisioning}~\cite{Zacarias2023Dynamic} starts the job with its request and every
five minutes resizes its allocation to the largest use of the coming five minutes, as the authors' simulator reads it from
the usage trace, local memory first and then the pool; a resize is committed only when it happens, and a resize that does
not fit fails the job, which restarts from the beginning. A running system cannot read the coming interval, so its
\term{look-back} variant resizes to the past five minutes' largest use, grows the job on demand a minute at a time, and
runs a job with its request held after three restarts, as the authors suggest for jobs that fail repeatedly. The
\term{Pond} oracle~\cite{Li2023Pond} holds the job's whole request but keeps only the request minus its predicted untouched
memory in local DRAM, the rest in the pool, with the untouched share predicted at the 0.98 quantile (Pond's target of 98\%
of jobs within their margin) and a QoS monitor that moves the pool part to local memory five minutes after the job first
touches it. The \term{KS+} oracle~\cite{KSplus2024} learns four segments of a non-decreasing step function from the RSS of
the program's other corpus inputs, regresses each segment's end and peak on the input size, ends segments 15\% earlier and
raises peaks 10\%, and after an OOM kill pulls the later segments forward to the failure or raises the last peak 20\%.

We checked these policies against their authors' published code. Our KS+ matches the authors' segmentation, prediction
and retry on random inputs. The dynamic-provisioning code of the authors' Slurm simulator matches ours in admission
against current allocations, remote-first release, local-first growth and restart with the original request; it also
recomputes the slowdown at every resize and charges 10\,s per growing resize, which we leave out in the policy's favour.
Pond's artifact holds none of its policies, so our Pond follows its paper.

Every estimate needs a margin, because an estimate below the truth causes OOM kills. A curve oracle's margin is headroom
added as a share of its own peak, and a peak oracle's margin multiplies its peak. Besides fixed margins of 10\% and 30\%,
we use \term{calibrated} margins that a site could set from its own history. For each oracle and program we compute the
smallest margin with which the oracle never falls below the truth, and a job receives the 90th percentile
(\term{calibrated 90\%}) or the largest value (\term{calibrated 100\%}) of the other programs' margins, so no oracle is
tuned on the program it predicts. The request, ML and reactive oracles have no profile to calibrate and keep a 10\% margin,
and the dynamic-provisioning, Pond and KS+ oracles use none.

\section{Footprint corpus}
\label{sec:corpus}

We ran 25 programs at three inputs each on a 32-core node with 62\,GB of memory (\tabref{tab:corpus}): four PolyBench
kernels~\cite{polybench}, miniVite~\cite{minivite}, the GAP benchmarks bfs, pr, cc and sssp, Lua, Python, Perl and SQLite
workloads, GCC's cc1 compiling three source files, HPCCG, LULESH, NPB IS, CG and MG, MiniFE, Quicksilver, Kripke, darknet
classifying images with AlexNet, and two LLM inference engines, llama2.c~\cite{llama2c} with three models and
llama.cpp~\cite{llamacpp} with SmolLM2-135M and 360M. Both engines map their weights from the model file rather than
allocating them. The native runs take 0.12 to 14.9\,s, and the true peaks at the largest inputs are 15 to 450\,MB. Every
(program, input) ran alone with its RSS sampled every 5\,ms; under a full instrumentation trace with Pin~\cite{Luk2005Pin},
which records the bytes touched over time and is the truth; under Massif and Massif with \texttt{--pages-as-heap},
each with its instruction clock and with milliseconds; and under DHAT, heaptrack and the gperftools heap profiler.

\begin{table}[t]
\centering
\caption{The corpus at each program's largest input: the true peak (bytes touched), the mean over the run as a share of the
peak, each method's peak as a share of the true peak, and the time each method takes over the native run.}
\label{tab:corpus}
\resizebox{\textwidth}{!}{\begin{tabular}{llrrrrrrrrrrr}
\toprule
 & & \multicolumn{6}{c}{Footprint (ratios to the true peak)} & & \multicolumn{4}{c}{Cost (time / native)} \\
\cmidrule(lr){3-8} \cmidrule(lr){10-13}
Program & Input & Peak (MB) & Mean/peak & RSS & Massif & DHAT & heaptrack & Native (s) & Massif & DHAT & heaptrack & Full trace \\
\midrule
2mm & LARGE & 37.1 & 0.96 & 1.03 & 1.00 & 1.00 & 1.00 & 8.6 & 5.6 & 15.2 & 2.3 & 357 \\
atax & LARGE & 32.1 & 0.87 & 1.04 & 1.00 & 1.00 & 1.00 & 0.1 & 5.7 & 9.8 & 2.5 & 122 \\
GCC cc1 & large & 105.5 & 0.69 & 1.27 & 0.16 & 0.15 & 0.15 & 1.8 & 11.0 & 68.1 & 2.2 & 206 \\
darknet & alexnet & 269.9 & 0.86 & 1.07 & 1.97 & 1.97 & 1.97 & 3.0 & 10.5 & 45.9 & 1.7 & 773 \\
GAP bfs & 18 & 72.8 & 0.36 & 1.05 & 0.99 & 0.99 & 0.99 & 1.5 & 3.0 & 5.7 & 1.6 & 83 \\
GAP cc & 18 & 72.0 & 0.36 & 1.06 & 1.01 & 1.01 & 1.01 & 2.2 & 2.2 & 3.9 & 1.1 & 56 \\
GAP pr & 18 & 72.4 & 0.40 & 1.05 & 1.00 & 1.00 & 1.00 & 2.8 & 2.3 & 5.3 & 1.2 & 123 \\
GAP sssp & 18 & 135.8 & 0.36 & 0.95 & 1.00 & 1.00 & 1.00 & 2.5 & 2.4 & 4.4 & 1.1 & 116 \\
gemm & LARGE & 29.1 & 1.00 & 1.04 & 1.00 & 1.00 & 1.00 & 6.1 & 5.7 & 15.2 & 1.1 & 356 \\
HPCCG & 64 & 102.7 & 0.99 & 1.03 & 1.02 & 1.02 & 1.02 & 1.7 & 10.5 & 24.8 & 1.4 & 310 \\
jacobi-2d & LARGE & 27.1 & 1.00 & 1.04 & 1.00 & 1.00 & 1.00 & 12.8 & 7.3 & 16.2 & 1.5 & 275 \\
Kripke & 16 & 274.6 & 0.88 & 1.02 & 1.00 & 1.00 & 1.00 & 3.4 & 4.8 & 50.7 & 1.4 & 473 \\
llama2.c & 110M & 449.7 & 0.98 & 1.14 & 0.17 & 0.17 & 0.17 & 14.9 & 3.6 & 20.5 & 1.3 & 41 \\
llama.cpp & 135m-n128 & 352.9 & 0.89 & 1.08 & 1.05 & 1.05 & 1.05 & 1.5 & 137.1 & 279.6 & 2.4 & 604 \\
Lua & 1 & 57.1 & 0.53 & 1.15 & 1.02 & 0.92 & 0.92 & 4.2 & 4.3 & 28.0 & 1.8 & 161 \\
LULESH & 60 & 151.1 & 0.62 & 1.02 & 1.00 & 1.00 & 1.00 & 5.5 & 7.9 & 32.2 & 1.3 & 458 \\
miniVite & 32768 & 15.4 & 0.51 & 1.99 & 1.55 & 1.53 & 1.42 & 1.8 & 22.9 & 36.9 & 4.5 & 322 \\
MiniFE & 80 & 256.5 & 0.98 & 1.04 & 1.03 & 1.02 & 1.02 & 9.2 & 6.7 & 23.8 & 1.1 & 273 \\
NPB CG & A & 27.4 & 1.00 & 1.10 & 0.99 & 0.99 & 0.99 & 1.9 & 7.8 & 15.5 & 1.2 & 237 \\
NPB IS & A & 69.7 & 0.81 & 1.03 & 0.99 & 0.99 & 0.99 & 2.4 & 2.7 & 7.0 & 1.0 & 61 \\
NPB MG & W & 58.5 & 0.98 & 1.04 & 0.99 & 0.99 & 0.99 & 0.5 & 4.0 & 8.8 & 1.3 & 107 \\
Perl & 1 & 56.1 & 0.67 & 1.15 & 1.04 & 0.97 & 0.96 & 1.8 & 5.2 & 44.8 & 1.9 & 145 \\
Python & 1 & 95.9 & 0.37 & 1.19 & 1.29 & 1.01 & 0.14 & 5.6 & 12.5 & 175.5 & 1.0 & 272 \\
Quicksilver & 100000 & 42.2 & 1.00 & 1.28 & 1.02 & 1.01 & 1.01 & 5.5 & 6.3 & 44.5 & 1.2 & 227 \\
SQLite & 1 & 77.8 & 0.43 & 1.06 & 1.01 & 1.01 & 1.01 & 1.2 & 7.1 & 39.0 & 1.6 & 204 \\
\bottomrule
\end{tabular}
}
\end{table}

\subsection{Accuracy of the peak}

\begin{figure}[t]
\centering
\includegraphics[width=\textwidth]{peak_accuracy.pdf}
\caption{Each method's peak over the true peak, every program at its largest input (one marker per program, the black bar
is the median). Left: against bytes touched. Right: against the RSS.}
\label{fig:peak}
\end{figure}

The heap profilers report the right peak for most of these programs and fail badly on the rest (\figref{fig:peak}).
Specifically, Massif's peak is within 5\% of the true peak on 20 of the 25 programs, DHAT and gperftools on 20 and
heaptrack on 19, because these programs touch nearly all the heap they allocate. However, the heap profilers read 0.14 to
0.16 times the peak of GCC's cc1, whose garbage collector maps its own pages, and 0.17 times that of llama2.c, whose
weights are mapped from the file and never pass through the heap; they read 1.97 times the peak of darknet and 1.42 to
1.55 times that of miniVite, which allocate memory they never touch. Massif with \texttt{--pages-as-heap} overestimates
every program (1.07 to 29.9 times), and the RSS is 0.95 to 1.99 times the bytes touched (median 1.05). A peak estimate is
therefore right or wrong by program, not by a small error spread over all programs.

\subsection{Shape over time}

\begin{figure}[t]
\centering
\includegraphics[width=0.6\textwidth]{motivation_curve.pdf}
\caption{Python (largest input): the true footprint over the run, the RSS, and the curves and peaks the heap profilers
report, each on its own method's clock rescaled to progress.}
\label{fig:motivation}
\end{figure}

The footprint of half of the programs varies over the run. The mean over the run is below 70\% of the peak for 11 of the
25 programs, and it is 0.36 to 0.40 times the peak for the four GAP kernels, so a scheduler that holds the peak for such a
job reserves memory that the job leaves unused for most of its run. \figref{fig:motivation} shows Python's footprint, which
rises and falls seven times. Massif's curve follows its phases on its instruction clock but sits above it, because Massif
counts requested heap bytes and Python allocates memory it never touches (1.29 times the true peak), and DHAT records only
the peak, its time and the heap at exit.

\begin{figure}[t]
\centering
\includegraphics[width=\textwidth]{llm_curves.pdf}
\caption{The two LLM engines at their largest inputs: the true footprint over the run, the RSS, and the curves and peaks
the heap profilers report.}
\label{fig:llmcurves}
\end{figure}

The LLM engines differ from the rest of the corpus in where their memory is (\figref{fig:llmcurves}). Both map the model's
weights from the file and read all of them for the first token, so the true footprint reaches 90\% of its peak within the
first 1 to 2\% of a llama2.c run and stays there, and llama.cpp reaches it at 12\% of a 128-token run. The heap profilers
see only the heap, 0.08 to 0.21 of llama2.c's peak, and the KV cache and compute buffers of llama.cpp.

\subsection{Cost}

\begin{figure}[t]
\centering
\includegraphics[width=0.6\textwidth]{cost.pdf}
\caption{Each method's time over the native run, every program at its largest input. The black bar is the median.}
\label{fig:cost}
\end{figure}

The exact footprint is expensive to obtain (\figref{fig:cost}). At the largest inputs the full instrumentation trace costs
226.8 times the native run at the median (41.2 to 772.5 times), against 1.4 times for heaptrack, 1.2 times for gperftools,
5.7 times for Massif (6.4 times with its clock in milliseconds) and 23.8 times for DHAT. On llama.cpp, a sub-second run
with four threads, the full trace costs 469 to 604 times the native run.

\section{Results}
\label{sec:results}

We first ask how much an exact footprint is worth and where the gain comes from, then how much of it each existing source
recovers and why each falls short. \tabref{tab:load08} gives every oracle at offered load 0.8 with calibrated 90\% margins,
and \figref{fig:load} gives the response time across loads.

\begin{table}[t]
\centering
\caption{Every oracle at offered load 0.8, calibrated 90\% margins (10\% for the request, ML and reactive oracles; none for
Pond, dynamic provisioning and KS+, whose Kills column counts failed resizes and retries), mean of three seeds. Resp.\ is
the mean response time, BSLD the mean bounded slowdown, Kills the OOM kills per stream of 3{,}000 jobs, Res.\ and Used the
memory reserved and resident as a share of all memory.}
\label{tab:load08}
\resizebox{\textwidth}{!}{\begin{tabular}{lrrrrrrrrrr}
\toprule
 & \multicolumn{5}{c}{Truth = bytes touched} & \multicolumn{5}{c}{Truth = RSS} \\
\cmidrule(lr){2-6} \cmidrule(lr){7-11}
Oracle & Resp. (min) & BSLD & Kills & Res. \% & Used \% & Resp. (min) & BSLD & Kills & Res. \% & Used \% \\
\midrule
Exact curve & 84 & 2.4 & 0 & 81 & 81 & 99 & 3.0 & 0 & 80 & 80 \\
Exact, 20 windows & 173 & 6.3 & 0 & 80 & 72 & 160 & 5.7 & 0 & 80 & 73 \\
Exact peak & 448 & 18.3 & 0 & 78 & 58 & 423 & 17.2 & 0 & 77 & 61 \\
Dynamic provisioning (5 min) & 457 & 18.9 & 307 & 74 & 63 & 395 & 16.1 & 265 & 75 & 64 \\
Dynamic provisioning, look-back & 504 & 21.1 & 444 & 73 & 61 & 493 & 20.5 & 478 & 72 & 61 \\
Reactive resizing (5 min) & 580 & 24.2 & 19 & 71 & 54 & 612 & 25.6 & 21 & 70 & 54 \\
RSS peak (earlier run) & 985 & 41.8 & 474 & 69 & 44 & 423 & 17.2 & 0 & 77 & 61 \\
DHAT peak & 1011 & 42.9 & 1323 & 70 & 45 & 1302 & 55.6 & 1310 & 65 & 41 \\
KS+ (segments from other inputs) & 1262 & 54.3 & 8189 & 59 & 38 & 844 & 35.7 & 8534 & 67 & 48 \\
ML peak (quantile GBM) & 1991 & 86.0 & 115 & 61 & 29 & 1992 & 86.1 & 115 & 61 & 31 \\
User request & 2208 & 95.6 & 0 & 57 & 27 & 2183 & 94.4 & 0 & 57 & 29 \\
Massif curve & 2399 & 103.8 & 1107 & 53 & 26 & 2099 & 90.6 & 1184 & 56 & 30 \\
Pond (untouched in pool) & 2519 & 109.0 & 0 & 53 & 26 & 2472 & 106.9 & 0 & 53 & 27 \\
DHAT curve (3 points) & 3512 & 152.9 & 1228 & 51 & 23 & 3008 & 130.7 & 1157 & 52 & 25 \\
DHAT hold-then-decay & 4161 & 180.9 & 1144 & 48 & 20 & 3943 & 171.4 & 1152 & 47 & 21 \\
Massif curve, ms clock & 4268 & 185.9 & 1056 & 45 & 17 & 2684 & 116.3 & 1190 & 51 & 26 \\
Massif pages-as-heap & 4670 & 203.7 & 668 & 40 & 15 & 4472 & 194.4 & 692 & 42 & 16 \\
Massif pages-as-heap, ms clock & 6619 & 290.0 & 465 & 35 & 10 & 5669 & 247.5 & 697 & 37 & 12 \\
gperftools peak & 10476 & 458.4 & 1644 & 33 & 7 & 11181 & 490.8 & 1644 & 33 & 7 \\
heaptrack curve & 11778 & 516.0 & 1285 & 31 & 6 & 12610 & 552.6 & 1439 & 31 & 6 \\
\bottomrule
\end{tabular}
}
\end{table}

\begin{figure}[t]
\centering
\includegraphics[width=\textwidth]{response_vs_load.pdf}
\caption{Mean response time against offered load, calibrated 90\% margins (95\% confidence intervals over three seeds, log
scale). Left: the truth is bytes touched. Right: the truth is the RSS.}
\label{fig:load}
\end{figure}

\subsection{What an exact footprint is worth}

Knowing a job's footprint curve, rather than its peak, is where most of the gain lies. At load 0.8 the exact curve gives a
mean response time of 84 minutes, against 448 minutes for the exact peak and 2{,}208 minutes for the user's request, under
the bytes-touched truth, which makes the curve 26.4$\times$ better than the request and 5.3$\times$ better than the peak,
and the peak 4.9$\times$ better than the request. The RSS truth gives the same picture (99, 423 and 2{,}183 minutes). The
curve lets the scheduler keep 81\% of all memory resident with jobs, against 58\% for the exact peak and 27\% for the
request, because jobs whose peaks fall at different times share the memory that a peak would hold idle. Coarsening the
exact curve to the 20 widened windows used for every measured curve raises the response to 173 minutes, so most of the
curve's advantage over the peak survives the resolution and timing slack a measured curve needs. The gap widens with load
(\figref{fig:load}): at load 0.4 every oracle that does not over-reserve keeps up with arrivals, and from load 0.6 the peak
and the request saturate memory while the curve does not.

\subsection{Existing footprint estimators}

No existing estimator comes close to the exact peak, let alone the exact curve. With calibrated 90\% margins, the RSS peak of
an earlier run gives 985 minutes and the DHAT peak 1{,}011, 2.2 times the exact peak, and the ML peak 1{,}991, little better
than the request it scales, because its 0.95-quantile prediction of $u$ covers only 76\% of jobs from unseen users. Massif's
curve gives 2{,}399 minutes, worse than the request, because the headroom that covers 90\% of the programs is 52\% of its
peak (\tabref{tab:margins}). Three kinds of program set that headroom: Massif reads GCC's cc1 at 0.16 times its peak and
llama2.c at 0.17, and on the four GAP kernels its clock of instructions drifts from the truth's clock enough to need 12 to
53\% headroom. With Massif's clock in milliseconds under Valgrind, which includes Valgrind's start-up and runs at a different
pace from the native program, the headroom rises to 100\% and the response time to 4{,}268 minutes. For gperftools and
heaptrack a third program falls above the 90th percentile, Python, whose heap they read at 0.21 and 0.14 times its peak, so
every job reserves three to four times its estimate and the response time rises to 10{,}476 and 11{,}778 minutes, about five
times the request's. The time axis of a curve is thus as important as its values, and an estimator that sees only the heap
is wrong by large factors on exactly the programs that set a site's margins.

\begin{table}[t]
\centering
\caption{The margin each oracle needs to stay above the truth, as a share of its own peak (\%), over the 25 programs at their
largest inputs.}
\label{tab:margins}
\begin{tabular}{lrrrrrr}
\toprule
 & \multicolumn{3}{c}{Truth = bytes touched} & \multicolumn{3}{c}{Truth = RSS} \\
\cmidrule(lr){2-4} \cmidrule(lr){5-7}
Oracle & Median & 90th pct. & Max & Median & 90th pct. & Max \\
\midrule
Massif curve & 0 & 52 & 534 & 2 & 56 & 689 \\
Massif pages-as-heap & 0 & 28 & 62 & 0 & 30 & 63 \\
Massif curve, ms clock & 0 & 100 & 535 & 20 & 71 & 690 \\
Massif pages-as-heap, ms clock & 0 & 65 & 78 & 0 & 52 & 76 \\
heaptrack curve & 1 & 326 & 593 & 2 & 376 & 747 \\
DHAT peak & 0 & 6 & 546 & 2 & 24 & 706 \\
DHAT curve (3 points) & 36 & 90 & 580 & 53 & 99 & 739 \\
DHAT hold-then-decay & 21 & 90 & 580 & 39 & 99 & 739 \\
gperftools peak & 0 & 233 & 554 & 2 & 290 & 716 \\
RSS peak (earlier run) & 0 & 1 & 7 & 0 & 0 & 0 \\
\bottomrule
\end{tabular}

\end{table}

\begin{figure}[t]
\centering
\includegraphics[width=\textwidth]{margin_tradeoff.pdf}
\caption{Mean response time against OOM kills at load 0.8, one marker per margin policy (10\%, 30\%, calibrated 90\% and
100\%). No margin gives an existing estimator both few kills and a short response time.}
\label{fig:margins}
\end{figure}

No margin policy rescues an estimator that is badly wrong for some programs (\figref{fig:margins}). A small margin lets those
programs' jobs be killed, and a margin large enough to cover them over-reserves every other job. With calibrated 100\%
margins every heap profiler must cover GCC's cc1, so each job reserves about 6.5 times its estimate and the response time
rises to 15{,}000 to 17{,}000 minutes. With a fixed 10\% margin, which accepts the kills, Massif gives 991 minutes with
1{,}714 kills per stream, and heaptrack and gperftools 1{,}154 and 1{,}158 minutes.

\subsection{Published provisioning policies}

The published policies recover part of the gain, each for a reason that a scheduler cannot rely on. \term{Dynamic
provisioning} gives 457 minutes under the bytes-touched truth, close to the exact peak (448), and 395 minutes under the RSS
truth, with 307 and 265 restarts per stream. It comes this close because each resize is the largest use of the coming five
minutes, read from the truth, so between resizes it is an exact curve at a five-minute resolution with no margin. Its
look-back variant, which knows only the past as a running system does, gives 504 and 493 minutes, behind the exact peak,
with 444 and 478 restarts and 25 and 27 jobs falling back to their request per stream. Reactive resizing, which also looks
back but reserves its resizes when the job starts, gives 580 and 612 minutes. A policy that resizes a running job thus
reaches the exact peak only with a view of the job's future, and pays for growth it does not foresee with restarts.

\term{Pond} gives 2{,}519 minutes, worse than the request (2{,}208), with no kills. It holds the whole request, so it cannot
pack more jobs than the request does, and the share it places in the pool must fit there even when the node's local memory
is free. Its model predicts 8\% of the request as untouched against 38\% in truth, and still overpredicts for 11\% of the
jobs, because they come from users the model has not seen. Pond saves local DRAM, not the total memory that the job streams
contend for.

\term{KS+} gives 1{,}262 minutes under the bytes-touched truth and 844 under the RSS truth, with about 8{,}200 and 8{,}500 OOM
kills per stream. Its first prediction of the largest input from the two smaller ones is 0.86 to 1.18 times the true RSS peak
for 19 of the 24 programs it can predict. Two programs cause nearly all the kills: the smaller inputs of atax and NPB MG
barely rise above the runtime's baseline, so the regressions predict almost nothing for the largest input, and their jobs
climb from the 100\,MB floor in 20\% steps, 36 to 38 times for atax and 20 to 23 times for MG. GCC's cc1 never fails but
reserves 8.6 to 10.7 times its peak. A model learned from a job's earlier executions with other inputs is thus only as good
as those inputs are close to the job's.

\subsection{LLM inference jobs}
\label{sec:llm}

\begin{figure}[t]
\centering
\includegraphics[width=\textwidth]{llm_response_vs_load.pdf}
\caption{A stream of LLM inference jobs alone (llama2.c with three models, llama.cpp with two lengths): mean response time
against offered load, calibrated 90\% margins, each configuration's margin calibrated on the other four.}
\label{fig:llm}
\end{figure}

\begin{table}[t]
\centering
\caption{The LLM stream at offered load 0.8, calibrated 90\% margins; columns as in \tabref{tab:load08}.}
\label{tab:llm}
\resizebox{\textwidth}{!}{\begin{tabular}{lrrrrrrrrrr}
\toprule
 & \multicolumn{5}{c}{Truth = bytes touched} & \multicolumn{5}{c}{Truth = RSS} \\
\cmidrule(lr){2-6} \cmidrule(lr){7-11}
Oracle & Resp. (min) & BSLD & Kills & Res. \% & Used \% & Resp. (min) & BSLD & Kills & Res. \% & Used \% \\
\midrule
Exact curve & 180 & 6.6 & 0 & 76 & 76 & 182 & 6.7 & 0 & 74 & 74 \\
Exact, 20 windows & 226 & 8.6 & 0 & 75 & 72 & 202 & 7.5 & 0 & 76 & 73 \\
Exact peak & 307 & 12.1 & 0 & 75 & 69 & 361 & 14.4 & 0 & 76 & 66 \\
Dynamic provisioning (5 min) & 386 & 15.8 & 140 & 72 & 65 & 443 & 18.3 & 200 & 71 & 63 \\
Dynamic provisioning, look-back & 528 & 22.2 & 260 & 68 & 62 & 694 & 29.3 & 356 & 62 & 56 \\
RSS peak (earlier run) & 598 & 24.9 & 0 & 69 & 57 & 361 & 14.4 & 0 & 76 & 66 \\
KS+ (segments from other inputs) & 737 & 31.1 & 1542 & 64 & 49 & 547 & 22.8 & 797 & 69 & 56 \\
Reactive resizing (5 min) & 753 & 31.9 & 10 & 65 & 54 & 712 & 29.9 & 9 & 66 & 54 \\
ML peak (quantile GBM) & 1836 & 79.1 & 115 & 61 & 36 & 1904 & 82.1 & 115 & 61 & 34 \\
User request & 2067 & 89.3 & 0 & 58 & 33 & 2146 & 92.7 & 0 & 57 & 31 \\
Pond (untouched in pool) & 2379 & 102.8 & 0 & 52 & 30 & 2450 & 105.9 & 0 & 53 & 28 \\
Massif pages-as-heap & 2650 & 114.7 & 0 & 47 & 26 & 2186 & 94.5 & 0 & 52 & 30 \\
Massif pages-as-heap, ms clock & 2955 & 127.9 & 0 & 45 & 24 & 2495 & 107.9 & 0 & 48 & 26 \\
Massif curve & 8747 & 381.2 & 1195 & 32 & 9 & 8528 & 372.3 & 1195 & 32 & 10 \\
heaptrack curve & 8794 & 383.2 & 1195 & 32 & 9 & 8333 & 363.2 & 1195 & 32 & 10 \\
Massif curve, ms clock & 8869 & 386.7 & 1195 & 32 & 9 & 8630 & 376.5 & 1195 & 32 & 10 \\
DHAT curve (3 points) & 9103 & 397.6 & 1195 & 31 & 9 & 8776 & 383.4 & 1195 & 32 & 9 \\
DHAT peak & 9198 & 400.7 & 1195 & 31 & 9 & 8986 & 391.7 & 1195 & 32 & 9 \\
gperftools peak & 9344 & 408.0 & 1195 & 31 & 9 & 9108 & 397.5 & 1195 & 32 & 9 \\
DHAT hold-then-decay & 9561 & 417.3 & 1195 & 31 & 8 & 9251 & 404.3 & 1195 & 32 & 9 \\
\bottomrule
\end{tabular}
}
\end{table}

On a stream of LLM jobs alone, whose memory is mostly read-only model weights, every heap profiler fails (\figref{fig:llm},
\tabref{tab:llm}). They read llama2.c at 0.08 to 0.21 of its peak, so the margin calibrated on the other configurations
cannot cover llama2.c's jobs: with calibrated margins Massif, DHAT, heaptrack and gperftools give 8{,}747 to 9{,}344 minutes,
four times the request (2{,}067). With a fixed 10\% margin, llama2.c's jobs are killed 9{,}050 to 15{,}661 times per stream
until the requeue factor has raised their reservation enough. The best existing source is the RSS peak of an earlier run,
598 minutes under the bytes-touched truth, because the weights the RSS counts are the footprint, against 180 minutes for the
exact curve and 307 for the exact peak. Dynamic provisioning gives 386 minutes by reading the coming five minutes, and its
look-back variant 528.

\subsection{A real trace}
\label{sec:lisa}

\begin{figure}[t]
\centering
\includegraphics[width=0.6\textwidth]{lisa_measured.pdf}
\caption{SURF Lisa jobs with their measured memory on a Lisa-shaped cluster, offered load 0.8: jobs per hour against local
memory per node (mean and 95\% interval over three seeds).}
\label{fig:lisa}
\end{figure}

\begin{table}[t]
\centering
\caption{The measured Lisa stream at offered load 0.8: jobs per hour and mean response time for each local memory per node
(mean of three seeds). The load is set against each cluster's memory, so numbers compare across rows, not across columns.}
\label{tab:lisa}
\resizebox{\textwidth}{!}{\begin{tabular}{lrrrrrrrrrrrr}
\toprule
 & \multicolumn{6}{c}{Jobs per hour} & \multicolumn{6}{c}{Mean response (min)} \\
\cmidrule(lr){2-7} \cmidrule(lr){8-13}
Local GB per node & 103 & 64 & 32 & 16 & 8 & 4 & 103 & 64 & 32 & 16 & 8 & 4 \\
\midrule
Exact curve & 9.6 & 9.4 & 8.6 & 3.5 & 3.8 & 2.0 & 11{,}790 & 10{,}755 & 9{,}735 & 36{,}491 & 36{,}939 & 14{,}073 \\
Exact, 20 windows & 9.6 & 9.4 & 7.5 & 3.2 & 3.5 & 1.9 & 11{,}790 & 10{,}755 & 12{,}264 & 46{,}250 & 46{,}383 & 20{,}655 \\
Exact peak & 9.6 & 9.4 & 6.6 & 2.7 & 2.9 & 1.7 & 11{,}790 & 10{,}755 & 16{,}080 & 59{,}143 & 60{,}184 & 28{,}962 \\
User request & 9.6 & 9.4 & 6.6 & 2.7 & 2.9 & 1.7 & 11{,}790 & 10{,}755 & 16{,}080 & 59{,}143 & 60{,}184 & 28{,}962 \\
Reactive resizing (5 min) & 9.6 & 9.4 & 7.5 & 3.0 & 3.4 & 1.8 & 11{,}790 & 10{,}755 & 11{,}705 & 46{,}413 & 44{,}847 & 18{,}800 \\
Dynamic provisioning (5 min) & 9.6 & 9.4 & 6.4 & 3.0 & 2.8 & 1.5 & 11{,}787 & 10{,}755 & 20{,}327 & 59{,}675 & 51{,}699 & 26{,}652 \\
Dynamic provisioning, look-back & 9.6 & 9.4 & 7.3 & 2.1 & 3.0 & 1.7 & 11{,}787 & 10{,}755 & 14{,}802 & 58{,}906 & 56{,}997 & 14{,}948 \\
\bottomrule
\end{tabular}
}
\end{table}

\begin{table}[t]
\centering
\caption{The hybrid Lisa stream (the trace's jobs with corpus footprint shapes) at 32\,GB per node and offered load 0.8,
calibrated 90\% margins. Fails counts OOM kills, restarts and retries.}
\label{tab:lisahybrid}
\resizebox{0.8\textwidth}{!}{\begin{tabular}{lrrrrrr}
\toprule
 & \multicolumn{3}{c}{Truth = bytes touched} & \multicolumn{3}{c}{Truth = RSS} \\
\cmidrule(lr){2-4} \cmidrule(lr){5-7}
Oracle & Jobs/h & Resp.\ (min) & Fails & Jobs/h & Resp.\ (min) & Fails \\
\midrule
Exact curve & 6.9 & 10{,}481 & 0 & 6.8 & 10{,}590 & 0 \\
Exact, 20 windows & 6.7 & 11{,}585 & 0 & 6.7 & 12{,}289 & 0 \\
Reactive resizing (5 min) & 6.7 & 11{,}740 & 2 & 6.5 & 12{,}983 & 4 \\
Dynamic provisioning, look-back & 6.6 & 13{,}451 & 251 & 6.6 & 13{,}749 & 282 \\
Dynamic provisioning (5 min) & 6.6 & 13{,}722 & 405 & 6.2 & 14{,}945 & 515 \\
User request & 6.3 & 15{,}563 & 0 & 6.3 & 16{,}012 & 0 \\
Exact peak & 6.3 & 15{,}563 & 0 & 6.3 & 16{,}012 & 0 \\
Pond (untouched in pool) & 6.0 & 19{,}710 & 0 & 6.0 & 19{,}677 & 0 \\
KS+ (segments from other inputs) & 4.8 & 19{,}920 & 6136 & 3.3 & 17{,}157 & 6411 \\
RSS peak (earlier run) & 5.8 & 20{,}623 & 474 & 6.3 & 16{,}012 & 0 \\
DHAT peak & 5.2 & 21{,}440 & 1323 & 4.9 & 24{,}514 & 1310 \\
ML peak (quantile GBM) & 4.1 & 21{,}892 & 212 & 4.0 & 21{,}682 & 212 \\
Massif curve & 3.9 & 31{,}021 & 1092 & 3.8 & 30{,}744 & 1188 \\
DHAT curve (3 points) & 3.1 & 45{,}306 & 1197 & 3.1 & 38{,}652 & 1149 \\
Massif curve, ms clock & 3.2 & 46{,}978 & 1040 & 3.6 & 34{,}698 & 1190 \\
heaptrack curve & 4.1 & 47{,}648 & 1285 & 3.8 & 45{,}731 & 1440 \\
gperftools peak & 3.3 & 59{,}035 & 1644 & 3.6 & 57{,}569 & 1644 \\
\bottomrule
\end{tabular}
}
\end{table}

On the Lisa trace with each job's measured memory (\figref{fig:lisa}, \tabref{tab:lisa}), knowing a job's memory matters once
memory is scarce. At Lisa's own 103\,GB and at 64\,GB per node every policy completes the same 9.4 to 9.6 jobs per hour, because
most jobs take whole 16-core nodes and use a few GB, so the cores are the bottleneck. From 32\,GB per node memory is
contended, and the exact curve completes 8.6 jobs per hour against 6.6 for the exact peak (30\% more, with a 40\% shorter
response time), and 18 to 31\% more from 16 down to 4\,GB. The peak and the request give the same results, because a
whole-node job blocks its node whatever memory it reserves. Dynamic provisioning restarts 775 to 2{,}129 jobs per stream at
32\,GB per node and below and gains less than reactive resizing. On the hybrid stream at 32\,GB per node
(\tabref{tab:lisahybrid}), the exact curve gives 10{,}481 minutes, reactive resizing 11{,}740, the two dynamic-provisioning
variants 13{,}451 and 13{,}722, the exact peak 15{,}563, Pond 19{,}710, KS+ 19{,}920 and the heap profilers 21{,}440 to
59{,}035 under the bytes-touched truth. On the trace's mix of job sizes, runtimes and node counts, the exact curve
therefore still leads, the resizing policies still fall between the curve and the peak, and the heap profilers still fall
behind the request, although the request now ties the exact peak because most jobs take whole nodes.

\subsection{Memory needed for the same service}

\begin{figure}[t]
\centering
\includegraphics[width=\textwidth]{capacity.pdf}
\caption{The total memory (32 nodes of local DRAM plus the 1\,TB pool) each oracle needs to keep the mean bounded slowdown
within 1.5 times that of the exact curve at 128\,GB per node, at load 0.8. Local memory per node was swept from 48 to
512\,GB.}
\label{fig:capacity}
\end{figure}

The response-time gains translate into memory a site would have to buy (\figref{fig:capacity}). At a fixed arrival rate, the
exact curve meets the target with 5.0\,TB in all. Under the bytes-touched truth the exact peak needs 6.0\,TB (20\% more),
reactive resizing and dynamic provisioning 7.0\,TB, DHAT, an earlier run's RSS, the ML peak and KS+ 9.0\,TB, Massif's curve
and the request 11.0\,TB (120\% more), the DHAT curve 13.0\,TB, and Massif on its millisecond clock and Pond 17.0\,TB.
gperftools and heaptrack do not meet the target even at 512\,GB per node.

\subsection{What each source's cost buys}

\begin{figure}[t]
\centering
\includegraphics[width=\textwidth]{roi.pdf}
\caption{Profiling cost per job (median time over the native run) against the gain in mean response time over the user's
request, at load 0.8, each oracle at its best margin policy. The exact curve's cost is the full trace under the
bytes-touched truth and an RSS sampler under the RSS truth.}
\label{fig:roi}
\end{figure}

The cost of a footprint and its value are far apart for every existing source (\figref{fig:roi}). Under the bytes-touched
truth the full trace costs 227 times the native run and gives a 26.4$\times$ gain over the request as a curve, or 4.9$\times$
as a peak. heaptrack and gperftools cost under 1.5 times the run and give 1.9$\times$, DHAT costs 24 times and gives
2.2$\times$, Massif's curve costs 5.7 times and gives 2.2$\times$ (1.3$\times$ with its millisecond clock), and an earlier
run's RSS costs nothing extra and gives 2.2$\times$. A source that is cheap gives a tenth of the exact curve's gain, and the
source that gives the exact curve costs two orders of magnitude more than the run.

\section{DHAT curves}
\label{sec:dhat}

DHAT records the heap at its global peak, the time of that peak, and the heap at exit, but no time
series~\cite{ValgrindDHAT}. A scheduler can use those three points to build a \term{DHAT curve}, which rises linearly from
zero to the peak at its recorded time and falls to the exit value, or a \term{hold-then-decay} curve, which holds the peak
until its recorded time and falls afterwards.

\begin{figure}[t]
\centering
\includegraphics[width=\textwidth]{dhat_curves.pdf}
\caption{DHAT's peak and the two curves built from its three recorded points, against the exact peak and curve, at load 0.8
under each margin policy (log scale).}
\label{fig:dhat}
\end{figure}

Both DHAT curves do worse than the DHAT peak under every margin policy and both truths, except with calibrated 100\% margins,
where all three reserve for GCC's cc1 (\figref{fig:dhat}). At load 0.8 with calibrated 90\% margins, the DHAT curve gives
3{,}512 minutes and the hold-then-decay curve 4{,}161, against 1{,}011 for the DHAT peak, under the bytes-touched truth. The
curves fall below the truth wherever a program has more than one phase, and the median program needs 36\% headroom on the
DHAT curve. The information that makes a curve useful to a scheduler is the timing of every phase, and DHAT records the
timing of one.

\section{Requirements for a footprint estimator}
\label{sec:req}

The results above determine what a source of footprint knowledge must provide for a scheduler to recover the gain of an
exact footprint.
\begin{itemize}
  \item \textbf{A curve, not a peak.} The exact curve is 5.3 times better than the exact peak at load 0.8 and needs 20\% less
    memory for the same service, so an estimator must report the footprint over the run.
  \item \textbf{A time axis that matches the run.} A curve on a clock other than the job's own (Massif's milliseconds under
    Valgrind, or its instructions on the GAP kernels) needs headroom that erases its advantage, and DHAT's single recorded
    time does not make a usable curve.
  \item \textbf{All of the memory the program touches, and only that.} The heap profilers miss memory a program maps itself
    (cc1's collector, the LLM engines' weights) and count memory it allocates and never touches (darknet, miniVite), and
    those programs set every site-wide margin.
  \item \textbf{A small margin that holds across programs.} An estimate that is exact for most programs and wrong by a large
    factor for a few forces a margin that over-reserves every job.
  \item \textbf{Known before the job starts.} Dynamic provisioning reaches the exact peak only by reading the job's future, and
    without it falls behind the peak and restarts hundreds of jobs per stream; Pond and KS+ need history that a new user or
    a new input does not have.
  \item \textbf{A cost close to the run's.} The full trace gives the exact curve at 227 times the native run, which no site can
    pay for every job.
\end{itemize}

\section{Limitations}
\label{sec:limits}

The results rest on a model of a cluster and of its jobs, and each simplification limits what they show.

\paragraph{The published baselines.} Our Pond follows its paper, because its artifact holds none of its policies, and it
leaves out the model that places latency-insensitive virtual machines entirely in the pool, which needs hardware counters our
jobs do not have. Our dynamic provisioning fixes a job's runtime stretch at admission and charges nothing per resize, both in
its favour. KS+ learns from two smaller inputs per program, where its authors had hundreds of executions per task.

\paragraph{Scaled footprints.} The corpus programs peak at 15 to 450\,MB, and the simulated jobs at tens to hundreds of GB.
Scaling a measured curve to a job's size keeps its shape and every estimator's ratio to the truth, but it assumes the ratio
does not change with the input, so we use the largest input of every program.

\paragraph{The job streams.} The main stream is synthetic, with memory the contended resource by design, and the
used-to-requested ratio comes from the LANL CM5 log of 1994 to 1996. The Lisa trace's jobs are those that had their nodes to
themselves, its memory is the nodes' active memory, and it records no memory request.

\paragraph{Scheduler and machine.} The slowdown is linear in the remote share of the footprint with one $\beta$ for every job,
the pool has no bandwidth contention, and only FCFS with EASY backfilling was tested. Appendix~\ref{sec:sensitivity} shows that the orderings
hold for $\beta$ from 0.1 to 2 and pools of 0.5 to 4\,TB, and that the update interval of the resizing policies changes their
results the most.

\paragraph{Costs.} The costs were measured with up to 15 runs sharing the node, and Pin's fixed start-up cost weighs heavily on
runs of a few seconds.

\section{Conclusion}
\label{sec:conclusion}

An exact footprint is worth most to a scheduler on disaggregated memory as a curve over the job's run, and the gain comes from
letting jobs with peaks at different times share memory: on our corpus and streams, the exact curve gives a 26.4$\times$
shorter mean response time than the user's request and needs 5.0\,TB where the request needs 11.0\,TB, and on the Lisa trace it
completes 18 to 31\% more jobs than the exact peak once memory is scarce. None of the existing sources recovers that gain. The
heap profilers are wrong by large factors on the programs that map their memory or never touch what they allocate, and no
margin covers those programs without over-reserving the rest. Of three published policies, dynamic provisioning comes close
to the exact peak only by reading each job's coming use, Pond saves local DRAM but not total memory, and KS+ cannot
extrapolate from inputs far from the job's. The exact curve itself costs 227 times the run. A footprint estimator for a
scheduler must therefore report a curve on the run's own time axis, count all the memory a program touches and nothing it
does not, hold with a small margin across programs, be known before the job starts, and cost close to the run. We present a
tool that targets these requirements in companion work~\cite{memprint}.

\appendix

\section{Sensitivity to the simulator's assumptions}
\label{sec:sensitivity}

\begin{table}[h]
\centering
\caption{Mean response time (minutes) at offered load 0.8 as the remote-memory slowdown $\beta$ varies (default 0.5),
calibrated 90\% margins, mean of three seeds.}
\resizebox{\textwidth}{!}{\begin{tabular}{lrrrrrrrrrr}
\toprule
 & \multicolumn{5}{c}{Truth = bytes touched} & \multicolumn{5}{c}{Truth = RSS} \\
\cmidrule(lr){2-6} \cmidrule(lr){7-11}
$\beta$ & 0.1 & 0.25 & 0.5 & 1 & 2 & 0.1 & 0.25 & 0.5 & 1 & 2 \\
\midrule
Exact curve & 58 & 63 & 84 & 190 & 390 & 58 & 67 & 99 & 201 & 470 \\
Exact peak & 274 & 340 & 448 & 693 & 1{,}162 & 233 & 302 & 423 & 666 & 1{,}194 \\
DHAT peak & 775 & 927 & 1{,}011 & 1{,}360 & 1{,}975 & 1{,}001 & 1{,}122 & 1{,}302 & 1{,}634 & 2{,}379 \\
User request & 1{,}896 & 2{,}012 & 2{,}208 & 2{,}601 & 3{,}365 & 1{,}846 & 1{,}977 & 2{,}183 & 2{,}627 & 3{,}431 \\
Reactive resizing (5 min) & 432 & 500 & 580 & 779 & 1{,}273 & 446 & 497 & 612 & 838 & 1{,}343 \\
Pond (untouched in pool) & 2{,}188 & 2{,}301 & 2{,}519 & 2{,}903 & 3{,}649 & 2{,}160 & 2{,}271 & 2{,}472 & 2{,}925 & 3{,}748 \\
Dynamic provisioning (5 min) & 286 & 332 & 457 & 568 & 1{,}021 & 274 & 363 & 395 & 594 & 932 \\
Dynamic provisioning, look-back & 481 & 439 & 504 & 686 & 1{,}121 & 460 & 507 & 493 & 706 & 1{,}064 \\
KS+ (segments from other inputs) & 1{,}031 & 963 & 1{,}262 & 1{,}378 & 1{,}828 & 704 & 741 & 844 & 1{,}215 & 1{,}282 \\
\bottomrule
\end{tabular}
}
\end{table}

\begin{table}[h]
\centering
\caption{Mean response time (minutes) at offered load 0.8 as the pool varies (default 1\,TB).}
\resizebox{\textwidth}{!}{\begin{tabular}{lrrrrrrrr}
\toprule
 & \multicolumn{4}{c}{Truth = bytes touched} & \multicolumn{4}{c}{Truth = RSS} \\
\cmidrule(lr){2-5} \cmidrule(lr){6-9}
Pool (GB) & 512 & 1024 & 2048 & 4096 & 512 & 1024 & 2048 & 4096 \\
\midrule
Exact curve & 565 & 84 & 51 & 51 & 671 & 99 & 51 & 51 \\
Exact peak & 1{,}499 & 448 & 78 & 52 & 1{,}455 & 423 & 65 & 52 \\
DHAT peak & 2{,}354 & 1{,}011 & 262 & 65 & 2{,}818 & 1{,}302 & 354 & 78 \\
User request & 4{,}441 & 2{,}208 & 697 & 261 & 4{,}458 & 2{,}183 & 667 & 222 \\
Reactive resizing (5 min) & 1{,}669 & 580 & 94 & 52 & 1{,}708 & 612 & 90 & 52 \\
Pond (untouched in pool) & 4{,}982 & 2{,}519 & 802 & 258 & 4{,}976 & 2{,}472 & 770 & 221 \\
Dynamic provisioning (5 min) & 1{,}624 & 457 & 58 & 50 & 1{,}617 & 395 & 57 & 50 \\
Dynamic provisioning, look-back & 1{,}887 & 504 & 56 & 50 & 2{,}134 & 493 & 56 & 50 \\
KS+ (segments from other inputs) & 2{,}005 & 1{,}262 & 437 & 126 & 1{,}669 & 844 & 350 & 113 \\
\bottomrule
\end{tabular}
}
\end{table}

\begin{table}[h]
\centering
\caption{Mean response time (minutes) at offered load 0.8 as the update interval of the policies that resize a running job
varies (default 5 minutes).}
\begin{tabular}{lrrrrrrrr}
\toprule
 & \multicolumn{4}{c}{Truth = bytes touched} & \multicolumn{4}{c}{Truth = RSS} \\
\cmidrule(lr){2-5} \cmidrule(lr){6-9}
Update interval (min) & 1 & 5 & 15 & 30 & 1 & 5 & 15 & 30 \\
\midrule
Reactive resizing & 321 & 580 & 1{,}034 & 1{,}402 & 329 & 612 & 1{,}063 & 1{,}421 \\
Dynamic provisioning & 157 & 457 & 812 & 1{,}330 & 182 & 395 & 872 & 1{,}232 \\
Dynamic provisioning, look-back & 182 & 504 & 1{,}175 & 1{,}657 & 191 & 493 & 1{,}029 & 1{,}758 \\
\bottomrule
\end{tabular}

\end{table}

We varied each of the simulator's assumptions alone, at offered load 0.8. The remote-memory slowdown $\beta$ changes every
oracle's response time, from 58 to 390 minutes for the exact curve between $\beta = 0.1$ and 2, but not their order. The pool
decides how much memory is contended: with 4\,TB most oracles come within a few minutes of the runtimes, and with 512\,GB the
differences grow (the exact curve 565 minutes, the exact peak 1{,}499, the request 4{,}441). The update interval matters most
for the policies that resize a running job: dynamic provisioning gives 157 minutes with an update every minute and 1{,}330
with one every 30 minutes under the bytes-touched truth, and its look-back variant 182 to 1{,}657 minutes.

\section{Reproducing the baselines' own results}
\label{sec:repro}

\begin{table}[h]
\centering
\caption{KS+'s memory wastage reduction against its baselines, as the KS+ paper reports it and as the authors' code gives it
on their traces, for 25, 50 and 75\% training data.}
\begin{tabular}{llrrrrrr}
\toprule
 & & \multicolumn{3}{c}{Paper} & \multicolumn{3}{c}{Reproduced} \\
\cmidrule(lr){3-5} \cmidrule(lr){6-8}
Workflow & KS+ against & 25\% & 50\% & 75\% & 25\% & 50\% & 75\% \\
\midrule
eager & k-Segments Selective & 36 & 39 & 40 & 36 & 40 & 40 \\
eager & PPM-Improved & 54 & 52 & 51 & 54 & 52 & 51 \\
sarek & k-Segments Selective & 31 & 28 & 29 & 32 & 28 & 30 \\
sarek & PPM-Improved & \multicolumn{3}{c}{45 (mean)} & 47 & 46 & 46 \\
\bottomrule
\end{tabular}

\end{table}

\begin{table}[h]
\centering
\caption{DRAM saved by pooling a fixed share of every virtual machine's memory, in percent, against the pool size in CPU
sockets, from the public Azure packing trace, with each server provisioned for its own peak or every server for the
cluster's largest peak; and Pond's Figure 3 at a 50\% share.}
\begin{tabular}{llrrrrrr}
\toprule
Provisioning & Pool share & 2 & 4 & 8 & 16 & 32 & 64 \\
\midrule
uniform fleet & 10\% & 3.8 & 3.9 & 4.0 & 4.0 & 4.1 & 4.1 \\
uniform fleet & 30\% & 11.5 & 11.8 & 11.9 & 12.1 & 12.2 & 12.3 \\
uniform fleet & 50\% & 19.2 & 19.6 & 19.9 & 20.2 & 20.4 & 20.6 \\
per server & 10\% & 0.0 & 0.1 & 0.2 & 0.3 & 0.4 & 0.4 \\
per server & 30\% & 0.0 & 0.4 & 0.7 & 1.0 & 1.2 & 1.3 \\
per server & 50\% & 0.0 & 0.6 & 1.2 & 1.6 & 1.9 & 2.2 \\
\midrule
Pond (paper) & 50\% & \multicolumn{4}{c}{} & 12 & 13 \\
\bottomrule
\end{tabular}

\end{table}

We reran each baseline's own evaluation where its data or a public substitute exists. The authors of KS+ publish their
simulation and traces, and every reduction their paper reports is reproduced within 2 percentage points. Zacarias et al.\
publish an artifact with their Slurm simulator, trace generation and a test set of nine simulations with output checksums;
built and run as its appendix prescribes, with its R packages taken from a CRAN snapshot of the artifact's date, all nine
outputs match the authors' checksums bit for bit. Pond's results rest on Azure's internal telemetry and its artifact holds
none of its policies, so only its Figure 3, the DRAM that pooling saves, can be rebuilt, from Azure's public packing trace.
Provisioning each server for its own peak reproduces the figure's shape (savings rising with pool size, more slowly beyond
16 sockets) at a sixth of its magnitude, and provisioning every server for the cluster's largest peak reproduces its magnitude
without its shape; the public trace lacks the server types and placements Pond used.

\bibliographystyle{plain}
\bibliography{references}

\end{document}
